\begin{section}{先进交互与虚拟现实实验室、北京触梦文化科技有限公司}
	\label{story:cucvr}

	2017年7月，张昊阳休学创立了\name{北京触梦文化科技有限公司}。

	2018年复学；7月，从社团事务中脱出身来，在学校创立\name{先进交互与虚拟现实实验室}。实验室位于国重15层1502，是一处风景极佳且没有宵禁的风水宝地。张昊阳将实验室开放给CUCVR社团内部活动用，一时间，1502成了核心社员们每夜party的固定去处。

	实验室帮扶校内创新创业大赛参赛者，屡获佳绩，故得以存续和运转，由信通学院综合实验中心提供场地和经费。

	同年秋，\name{游戏研发社}时任社长\name{王思遥}因无处开展社团活动	人员稀少濒临解散，找到张昊阳寻求帮助，遂亦将1502开放给游研社。

	年末，张昊阳作为北京市文创领域大学生创业者代表，参加北京市市委宣传部组织的集中培训，期间认识了某文化传媒公司（疑为\name{北京锦熹文化传媒有限公司}）。

	2019年4月，实验室若干成员通过触梦公司中标了为上述文化传媒公司提供数字艺术制作服务的机会。

	是夏，实验室核心成员有：张昊阳、赵一桥、吕鸿风、沈琰周、王思遥、王韬涵、周子腾、王思萌。

	项目最终流标，触梦公司起诉该文化公司。

	2020年暑假，实验室被前来“参观”的老师设法抢走。

	2021年，触梦公司最终败诉，张昊阳身负数十万元债务。
\end{section}